\documentclass[10pt,a4paper]{amsart}
\usepackage[margin=19mm]{geometry}
\usepackage{amsmath,amssymb,mathtools}
\usepackage[scr=rsfs,cal=euler]{mathalfa}
\usepackage{xfrac}
\usepackage[unicode,hidelinks,bookmarks=false]{hyperref}
\newcommand{\ie}{i.e.\ }
\newcommand{\cf}{cf.\ }
\newcommand{\resp}{respectively\ }
\newcommand{\Subsection}{\subsection}
\providecommand{\nobreakdash}{\nobreak\hbox{-}}
\setlength{\emergencystretch}{2em}
\multlinegap0pt
\usepackage{amssymb}
\usepackage{float}
\usepackage{mathtools}
\usepackage{cancel}
\usepackage{tcolorbox}
\usepackage{bm}
\usepackage{csquotes}
\usepackage[shortlabels]{enumitem}
\usepackage{xcolor}
\newtheorem{theorem}{Theorem}
\title{Assembly Addition Chains}
\author{Leroy Cronin, Juan Carlos Morales Parra, Keith Y. Patarroyo}
\date{Source: arXiv:2512.18030v1 (CC BY 4.0)}
\begin{document}
\maketitle

\noindent\emph{Source and licence.} Assembly Addition Chains. Version arXiv:2512.18030v1 (CC BY 4.0), distributed by arXiv under the Creative Commons Attribution 4.0 International licence, http://creativecommons.org/licenses/by/4.0/, as recorded in the arXiv licence metadata for this exact version and shown on the abstract page https://arxiv.org/abs/2512.18030v1. Full licence notice: \texttt{licenses/arxiv-2512.18030-CC-BY-4.0.txt}.\par\smallskip

\noindent\textit{Editorial scope.} Counted unit: the article's theorem ineq2PD (source0.tex lines 1501--1506). The article's complete original proof of it is delivered with it (source0.tex lines 1507--1533, one \texttt{proof} environment of 27 lines). No mathematics is written, repaired or re-derived: the statement and the proof are the article's own lines, in the article's own notation, and no retained line is substituted or re-flowed.  Retained uncounted as the source-local context the proof needs: lines 1477--1479; lines 1485--1487.  Nothing was omitted from any retained span, so the package carries no omission record.  No retained line contains a citation, so the wrapper carries no bibliography and no bibliography entry.  No retained line contains a figure environment, an image-inclusion command or drawing code, so no image is delivered and the package declares no asset dependency.  Editorial formatting only, in addition: an AMS \texttt{amsart} preamble that restates the article's own declarations and macros used by the retained lines, namely the environments \texttt{theorem} on separate counters, together with the standard packages those lines need.  The retained environments print in this document as Theorem 1 in order of appearance, whereas the article prints the same items with the numbering of its own full document.  The complete native source of the article as distributed in the arXiv e-print for this exact version is bundled unchanged as \texttt{sources/191356/source0.tex}.
        \begin{equation}\label{2piece2k}
        \mathscr{O} \in ((\mathscr{B}_{\sigma(1)} \circ \mathscr{B}_{\sigma(2)}) \circ \cdots \circ ((\mathscr{B}_{\sigma(2k-3)} \circ \mathscr{B}_{\sigma(2k-2)}) \circ (\mathscr{B}_{\sigma(2k-1)} \circ \mathscr{B}_{\sigma(2k)}\underbrace{)) \cdots )}_{k}.
        \end{equation}

        \begin{equation}\label{2piece2k1}
            \mathscr{O} \in ((\mathscr{B}_{\sigma(1)} \circ \mathscr{B}_{\sigma(2)}) \circ \cdots \circ ((\mathscr{B}_{\sigma(2k-3)} \circ \mathscr{B}_{\sigma(2k-2)}) \circ (\mathscr{B}_{\sigma(2k-1)} \circ \mathscr{B}_{\sigma(2k)}\underbrace{)) \cdots )}_{k} \circ \mathscr{B}_{\sigma(2k+1)}.   
        \end{equation}

\begin{theorem}
Let $(S,\circ,BB)$ be an Assembly Space. For any $\mathscr{O} \in 2PD(S)$ we have
\begin{equation}\label{ineq2PD}
    a(\mathscr{O}) \leq Ma_{2PD}(\mathscr{O}) \leq s(\mathscr{O})-1.
\end{equation}
\end{theorem}

\begin{proof}
Let $\mathscr{O} \in 2PD(S)$. If $s(\mathscr{O})=2k$ for $k \in \mathbb{Z}^+$, then by definition there exists $C \in AAC(\mathscr{O},BB)$ given in terms of (\ref{2piece2k}) by  
\begin{equation*}
\begin{split}
   C= & (\mathscr{B}_{\sigma(1)} \circ \mathscr{B}_{\sigma(2)}, \dots,\mathscr{B}_{\sigma(2k-1)} \circ \mathscr{B}_{\sigma(2k)}, (\mathscr{B}_{\sigma(2k-3)} \circ \mathscr{B}_{\sigma(2k-2)}) \circ (\mathscr{B}_{\sigma(2k-1)} \circ \mathscr{B}_{\sigma(2k)}),\dots,\\& ((\mathscr{B}_{\sigma(1)} \circ \mathscr{B}_{\sigma(2)}) \circ \cdots \circ ((\mathscr{B}_{\sigma(2k-3)} \circ \mathscr{B}_{\sigma(2k-2)}) \circ (\mathscr{B}_{\sigma(2k-1)} \circ \mathscr{B}_{\sigma(2k)}\underbrace{)) \cdots )}_{k}). 
   \end{split}
\end{equation*}
Since the first $k$ objects in the sequence have size two and the number of objects of size two in $S$ is $\text{Card}(S(2))$, then there exists $C' \in AAC(\mathscr{O},BB)$ (obtained from $C$ by removing possible repeated objects of size two) such that  
\begin{equation*}
    L(C')= \min \{k, \text{Card}(S(2)) \} + k-1.
\end{equation*}
Analogously, if $s(\mathscr{O})=2k+1$ for some $k \in \mathbb{Z}^+$, then there exists $C \in AAC(\mathscr{O},BB)$ given implicitly by (\ref{2piece2k1})
\begin{equation*}
\begin{split}
   C= & (\mathscr{B}_{\sigma(1)} \circ \mathscr{B}_{\sigma(2)}, \dots,\mathscr{B}_{\sigma(2k-1)} \circ \mathscr{B}_{\sigma(2k)}, (\mathscr{B}_{\sigma(2k-3)} \circ \mathscr{B}_{\sigma(2k-2)}) \circ (\mathscr{B}_{\sigma(2k-1)} \circ \mathscr{B}_{\sigma(2k)}),\dots,\\& ((\mathscr{B}_{\sigma(1)} \circ \mathscr{B}_{\sigma(2)}) \circ \cdots \circ ((\mathscr{B}_{\sigma(2k-3)} \circ \mathscr{B}_{\sigma(2k-2)}) \circ (\mathscr{B}_{\sigma(2k-1)} \circ \mathscr{B}_{\sigma(2k)}\underbrace{)) \cdots )}_{k},\\
   & ((\mathscr{B}_{\sigma(1)} \circ \mathscr{B}_{\sigma(2)}) \circ \cdots \circ ((\mathscr{B}_{\sigma(2k-3)} \circ \mathscr{B}_{\sigma(2k-2)}) \circ (\mathscr{B}_{\sigma(2k-1)} \circ \mathscr{B}_{\sigma(2k)}\underbrace{)) \cdots )}_{k}\circ \mathscr{B}_{2k+1}) 
   \end{split}
\end{equation*}
So, there exists $C' \in AAC(\mathscr{O},BB)$ (obtained from $C$ by removing possible repeated objects of size two) such that 
\begin{equation*}
    L(C')= \min \{k, \text{Card}(S(2)) \} + k.
\end{equation*}
The last inequality in (\ref{ineq2PD}) simply follows from 
\begin{equation*}
   \bigg\lfloor \frac{s(\mathscr{O}}{2} \bigg\rfloor + \bigg\lceil \frac{s(\mathscr{O})}{2} \bigg\rceil = s(\mathscr{O}).
\end{equation*}
\end{proof}

\end{document}
